% Enterprise AI Safeguarding Survey Draft
\documentclass[sigconf]{acmart}

% Draft metadata: copyright, DOI, and conference parameters configured for preprint/review.
\settopmatter{printacmref=false}
\renewcommand\footnotetextcopyrightpermission[1]{}

\begin{document}

\title{Safeguarding Autonomous Enterprise AI Agents: A Survey of Access Control, Threat Mitigation, and Governed Execution (2020--2026)}

\author{[Author Name(s) -- replace before submission]}
\affiliation{%
  \institution{[Department / Institution -- replace]}
  \city{[City -- replace]}
  \country{[Country -- replace]}
}
\email{[Corresponding author email -- replace]}
% Duplicate the author/affiliation/email block above for each additional author.

\begin{abstract}
Enterprise AI is rapidly transitioning from passive conversational assistants to autonomous agentic systems. These agents connect large language models to sensitive corporate data stores and real-world execution APIs, introducing critical vulnerabilities that extend far beyond inaccurate text generation. An unconstrained agent can expose confidential documents across tenant boundaries, execute malicious directives embedded in retrieved content, or trigger destructive, non-atomic tool mutations. This survey provides a structured review of academic research (2020--2026) and industrial practices for safeguarding enterprise AI agents. We synthesize existing literature into a four-layer taxonomy spanning the complete data-to-action execution lifecycle: (1) permission-aware data governance and grounding, (2) threat mitigation and prompt-injection defense, (3) governed tool execution and action reliability, and (4) operational observability and lifecycle auditing. We evaluate model-based guardrails, native access-control integration, capability-based information flow control, formal sequence verification, and postcondition-aware idempotency ledgers. Grounded by an empirical case study of an industrial multi-agent platform, our analysis reveals that reliable safety cannot be achieved through prompt engineering or model alignment alone. Rather, enterprise safeguarding is fundamentally a systems-engineering challenge requiring deterministic, software-enforced boundaries across each stage of agent execution. Finally, we highlight critical research gaps in synchronization, production evaluation fidelity, and unified forensic auditability.
\end{abstract}

\keywords{enterprise AI, autonomous agents, retrieval-augmented generation, access control, prompt injection, AI guardrails, tool safety, idempotency}

\maketitle
\pagestyle{plain} % Avoid template placeholder conference headers in this draft.

% ==================================================================
% 1. INTRODUCTION
% ==================================================================
\section{Introduction}
Modern enterprises are increasingly deploying autonomous AI agents to orchestrate multi-stage business workflows \cite{yao2023react,wang2024agents}. Unlike traditional, passive question-answering systems, enterprise agents actively ingest corporate documentation from relational and vector stores, synthesize contextual plans, and execute consequential business actions across internal and external APIs (e.g., dispatching bulk emails, modifying customer records, or invoking financial transactions). In this agentic paradigm, an error is no longer confined to inaccurate text generation; a single failure can trigger unauthorized data disclosure, severe operational disruption, or irreversible side effects.

Securing these systems introduces a multidimensional challenge across the data-to-action pipeline:
\begin{enumerate}
  \item \textbf{Data Leakage:} A user querying the corporate knowledge base may receive retrieved context containing proprietary or confidential records belonging to other tenants or teams.
  \item \textbf{Adversarial Manipulation:} A retrieved business document or external webpage may contain embedded adversarial instructions that hijack the model's control flow.
  \item \textbf{Unreliable Tool Execution:} An authorized tool execution may fail ambiguously due to network latency, causing naive retry logic to duplicate real-world actions.
  \item \textbf{Financial and Operational Exhaustion:} An agent trapped in recursive tool loops can incur runaway API costs or trigger unmonitored system mutations.
\end{enumerate}

These operational hazards cannot be resolved through an isolated output filter or system-prompt instructions. As codified in authoritative taxonomies such as the OWASP Top 10 for Large Language Models \cite{owasp2025llm} and structured AI threat modeling frameworks \cite{dev2025threat}, safeguarding enterprise agents demands explicit, architectural mechanisms. Furthermore, embedding operational and ethical guardrails requires systematic integration across the software development lifecycle (SDLC) rather than relying on probabilistic model self-governance \cite{kulkarni2025grails,dong2024systematicguardrails}.

This survey synthesizes peer-reviewed literature, security standards, and industrial engineering practices published between 2020 and 2026. To systematically analyze the state of the art and emerging engineering practices, our review is guided by four foundational Research Questions (RQs) that directly correspond to each stage of the agent execution lifecycle:
\begin{itemize}
  \item \textbf{RQ1 (Data Governance \& Permissions):} \emph{How can enterprise AI agents retrieve contextually relevant corporate knowledge while enforcing multi-tenant isolation and source-specific access control policies at query time?}
  \item \textbf{RQ2 (Threat Protection \& Injection Defense):} \emph{How can autonomous agents mitigate Indirect Prompt Injection (IPI) and prevent untrusted external data from hijacking model control flows or sensitive execution parameters?}
  \item \textbf{RQ3 (Governed Tool Execution \& Action Reliability):} \emph{How can tool-mediated agent actions be constrained to authorized, validated operations, and how can state consistency and idempotency be guaranteed under non-atomic execution failures?}
  \item \textbf{RQ4 (Operational Observability \& Lifecycle Auditing):} \emph{How should multi-layered enterprise safeguards be systematically evaluated, continuously monitored, and forensically audited across the software development lifecycle?}
\end{itemize}

The remainder of this paper is structured as follows: Section~\ref{sec:taxonomy} establishes foundational preliminaries and introduces our four-layer safeguard taxonomy; Section~\ref{sec:review} delivers a critical thematic review addressing each of the four RQs; Section~\ref{sec:comparative} provides a comparative analysis and examines an empirical codebase case study; Section~\ref{sec:gaps} identifies key open research gaps; and Section~\ref{sec:conclusion} concludes with core takeaways.

% ==================================================================
% 2. BACKGROUND AND TAXONOMY
% ==================================================================
\section{Background and Taxonomy}
\label{sec:taxonomy}

\subsection{Preliminaries: Enterprise Agent Workflows}
An enterprise agent system operates over an iterative data-to-action execution loop:
\[
\text{Data Ingestion} \longrightarrow \text{Retrieval \& Grounding} \longrightarrow \text{Model Reasoning} \longrightarrow \text{Tool Execution} \longrightarrow \text{State Mutation}
\]
Knowledge retrieval is commonly implemented via Retrieval-Augmented Generation (RAG) \cite{gao2023ragsurvey}, which projects corporate documents into dense vector spaces and retrieves candidate chunks using similarity measures (e.g., cosine distance). Agentic reasoning follows frameworks such as ReAct (Reasoning and Acting) \cite{yao2023react}, where the language model interleaves chain-of-thought planning with structured tool calls. Because tools interact directly with production databases and external communication services, the execution environment must enforce strict operational constraints.

\subsection{The Four-Layer Enterprise Safeguard Taxonomy}
To organize the diverse safeguards proposed across the literature, we classify them into a four-layer architectural taxonomy directly mapped to our research questions:
\begin{itemize}
  \item \textbf{Layer 1 (RQ1): Data Governance \& Permission-Aware Grounding (What may the agent know?)} Controls that enforce identity, multi-tenant isolation, and source-specific access policies before data enters the model context, alongside metrics that evaluate factual faithfulness.
  \item \textbf{Layer 2 (RQ2): Threat Protection \& Injection Defense (How is untrusted input constrained?)} Guardrails and structural isolation techniques that detect adversarial prompts, filter sensitive entities, and prevent untrusted content from altering execution flows.
  \item \textbf{Layer 3 (RQ3): Governed Tool Execution \& Action Reliability (What may the agent do?)} State machines, formal sequence constraints, human approval gates, and verify-before-retry ledgers that govern external side effects and enforce idempotency.
  \item \textbf{Layer 4 (RQ4): Operational Observability \& Lifecycle Auditing (How is the system monitored?)} Frameworks for AI threat modeling, software lifecycle governance, token budget caps, and tamper-resistant telemetry.
\end{itemize}

\subsection{Review Methodology}
We conduct a thematic literature review spanning peer-reviewed conference proceedings, journal articles, and technical reports published between 2020 and 2026 across major computing and security venues (ACM, IEEE, ACL, USENIX, arXiv). Candidate works are selected based on their applicability to enterprise RAG and multi-agent governance, and are analyzed critically according to their threat models, formal guarantees, and empirical validation.

% ==================================================================
% 3. THEMATIC LITERATURE REVIEW
% ==================================================================
\section{Thematic Literature Review}
\label{sec:review}

\subsection{RQ1: Data Governance, Permissions, and Grounding}
A standard RAG pipeline optimizes for semantic relevance rather than authorization. When enterprise documents originating from heterogeneous services (e.g., AWS S3, Google Drive, Jira) are indexed into a unified vector store, vector similarity search operates oblivious to the querying user's access rights. Attempting to translate and re-implement diverse enterprise permissions into an ad-hoc centralized database frequently introduces mapping flaws that expand access.

Jeong and Lee's \emph{Permission-Aware RAG} \cite{jeonglee2025permissionawarerag} addresses this authorization gap by decoupling semantic retrieval from access control. Their architecture first retrieves candidate context chunks based on vector proximity and then queries the native Identity and Access Management (IAM) engine of the originating service (e.g., AWS IAM, GCP IAM, Keycloak) before chunks are admitted into the prompt context. By evaluating native permissions in parallel at query time, the system avoids error-prone policy synchronization. In experiments on HotpotQA with 13 simulated user profiles, native validation achieved 100\% allow/deny agreement, whereas a merged-policy baseline misclassified access due to missing attribute conditions \cite{jeonglee2025permissionawarerag}. However, permission enforcement exhibits a distinct latency trade-off: scaling top-$k$ from 1 to 20 increased mean IAM validation latency from 65\,ms to 204\,ms, directly impacting interactive responsiveness.

Even when retrieved context is legally authorized, language models can generate unfaithful or hallucinated claims. Es et al.'s \emph{RAGAs} framework \cite{es2024ragas} formalizes automated, reference-free evaluation for RAG pipelines. RAGAs decomposes response quality into three independent dimensions: \emph{context relevance} (filtering noise from retrieved chunks), \emph{faithfulness} (verifying that generated claims are mathematically grounded in the context), and \emph{answer relevance} (measuring query alignment). While RAGAs effectively identifies unsupported parametric generation, it operates as an informational metric; it cannot prevent unauthorized context ingestion or adversarial manipulation.

\subsection{RQ2: Threat Protection and Prompt Injection Defense}
Integrating external data sources exposes agents to severe threat vectors, most notably Indirect Prompt Injection (IPI). As first formalized by Greshake et al. \cite{greshake2023indirect}, IPI occurs when an untrusted third-party document (e.g., a customer email or vendor invoice) embeds adversarial instructions designed to subvert the agent's task, exfiltrate private context, or trigger unauthorized API calls.

Existing literature approaches threat mitigation through two distinct paradigms:
\begin{enumerate}
  \item \textbf{Model-Based Guardrails and Classifiers:} Dong et al.'s broad survey \cite{dong2025safeguardingsurvey} and position paper \cite{dong2024systematicguardrails} categorize input/output guardrails into safety classifiers, pattern-matching rules, and alignment fine-tuning. Meta's LlamaGuard \cite{inan2023llamaguard} provides a standardized classification baseline across safety categories. Safe-BigQA \cite{moraes2026safebigqa} optimizes this model-based approach for enterprise systems by deploying a two-stage filter: a lightweight, low-latency classifier screens incoming prompts, routing borderline inputs to a secondary LLM guard. The authors report jailbreak-detection F1 scores up to 62.87\%, reflecting a 9.64\% gain over standalone LlamaGuard \cite{moraes2026safebigqa}. Similarly, OpenGuardrails \cite{wang2025openguardrails} offers modular entity redaction, while Protect \cite{avinash2025protect} extends guardrail stacks to multimodal modalities via LoRA adapters. SANS practitioner evaluations \cite{mcwherter2024safeguardingai} observe that combining system instructions with NVIDIA NeMo self-checks blocks common jailbreaks where prompts alone fail.
  \item \textbf{Capability-Based Information Flow Control:} While model classifiers improve hygiene, they are probabilistic and vulnerable to novel adversarial circumventions. Debenedetti et al.'s \emph{CaMeL} \cite{debenedetti2026camel} establishes a structural defense against IPI by enforcing strict privilege separation. CaMeL wraps the agent in an architectural execution layer that extracts data-flow and control-flow graphs directly from trusted user instructions. Untrusted external context is processed exclusively by a tool-less reader model; data originating from untrusted sources is tagged and cryptographically constrained by capabilities, preventing it from ever influencing sensitive tool arguments (such as an outbound email recipient). On the AgentDojo benchmark, CaMeL solved 77\% of tasks with formal security guarantees, compared to an 84\% baseline for undefended agents \cite{debenedetti2026camel}.
\end{enumerate}

\subsection{RQ3: Governed Tool Execution and Action Reliability}
Granting an agent access to external tools introduces the risk of excessive agency (OWASP LLM08 \cite{owasp2025llm}), where an autonomous model executes actions beyond its intended scope. Reliable action governance requires enforcing deterministic software boundaries between model intention and execution.

Doshi et al. \cite{doshi2026verifiablysafetooluse} propose a formal method for tool safety by applying System-Theoretic Process Analysis (STPA) to enterprise workflows. By systematically identifying hazards (e.g., sending unauthorized data or deleting records), they express allowable tool sequences and information flows in Alloy, proving safety invariants via formal model checking. 

Complementary to authorization is the challenge of \emph{action reliability}. In production environments, tool execution is rarely atomic. Network partitions, timeouts, and delayed visibility frequently return ambiguous errors. If an agent naively retries an ambiguous tool invocation, it risks repeating non-idempotent operations (such as charging a credit card or dispatching duplicate emails). Mansoor et al.'s \emph{Verified Tool Calls} \cite{mansoor2026verifiedtoolcalls} addresses non-atomic failures by wrapping tool executions in a verification-aware protocol. The system mandates verify-before-retry logic: following an ambiguous failure, the agent queries external system state (postconditions) and supplies idempotency keys, eliminating duplicate actions across tested failure scenarios \cite{mansoor2026verifiedtoolcalls}.

\subsection{RQ4: Operational Observability and Lifecycle Auditing}
Operational governance requires continuous monitoring of agent behavior, resource expenditure, and safety compliance. Dev et al. \cite{dev2025threat} conduct a meta-analysis of 14 AI threat modeling frameworks (including MITRE ATLAS, OWASP, and NIST AI RMF), demonstrating that guardrail implementations frequently fail because they lack structured threat definitions. They propose a standardized threat library allowing practitioners to map specific application risks to concrete guardrail configurations.

To operationalize these safeguards throughout enterprise engineering, Kulkarni and Ramanathan introduce the \emph{GRAILS} framework \cite{kulkarni2025grails}. GRAILS establishes architectural design patterns for embedding ethical, legal, and operational guardrails into the software development lifecycle (SDLC). Rather than treating safety as an external post-hoc wrapper, GRAILS integrates policy enforcement, budget reservations, and audit checkpoints directly into software specifications, ensuring verifiable governance from design to production.

% ==================================================================
% 4. COMPARATIVE ANALYSIS
% ==================================================================
\section{Comparative Analysis}
\label{sec:comparative}

\subsection{Comparison with Existing Literature and Related Surveys}
A key limitation of existing research on LLM safety is fragmentation: prior surveys and primary research papers typically isolate a single layer—such as conversational output guardrails, vector retrieval algorithms, or isolated injection benchmarks—without addressing the interconnected failure modes of multi-stage autonomous agents. 

To clarify the scope and contributions of existing literature, Table~\ref{tab:survey_comparison} compares this survey against prior surveys, meta-reviews, and representative primary research papers across five key architectural dimensions:
\begin{enumerate}
  \item \textbf{Primary Focus:} The principal problem domain addressed by the work (e.g., conversational text safety, RAG retrieval mechanics, autonomous agent design, or isolated tool safety).
  \item \textbf{Data Access \& IAM (RQ1):} Whether the work enforces multi-tenant identity and access management (IAM) to prevent unauthorized document exposure during retrieval.
  \item \textbf{Prompt Injection Defense (RQ2):} Whether the approach provides structural or classifier-based defenses against indirect prompt injection embedded in external data.
  \item \textbf{Tool \& Action Safety (RQ3):} Whether the system enforces runtime constraints, idempotency ledgers, or human approval checkpoints before executing real-world external mutations.
  \item \textbf{Standards Alignment (RQ4):} Whether the work formally aligns with recognized enterprise security frameworks, such as the OWASP Top 10 for LLMs \cite{owasp2025llm}, NIST AI RMF, or MITRE ATLAS \cite{dev2025threat}.
\end{enumerate}

As shown in Table~\ref{tab:survey_comparison}, existing broad surveys on LLM safeguards (e.g., Dong et al.~\cite{dong2025safeguardingsurvey}) and RAG (Gao et al.~\cite{gao2023ragsurvey}) focus primarily on conversational text safety and semantic retrieval, lacking operational tool governance or multi-tenant IAM enforcement. General surveys on autonomous agents (e.g., Wang et al.~\cite{wang2024agents}) analyze planning and tool selection but omit security boundaries, IAM controls, and prompt-injection defenses. Conversely, specialized research papers (such as Jeong and Lee~\cite{jeonglee2025permissionawarerag} for IAM-RAG, Debenedetti et al.~\cite{debenedetti2026camel} for CaMeL, or Mansoor et al.~\cite{mansoor2026verifiedtoolcalls} for tool idempotency) offer deep mechanisms for individual boundaries but do not cover the complete data-to-action pipeline. This survey bridges these isolated silos by delivering a unified, four-layer systems taxonomy (RQ1--RQ4) grounded in both foundational security standards and real-world system architecture.

\begin{table*}[t]
\centering
\caption{Comparison of our survey with existing surveys, meta-reviews, and representative related works on AI safeguarding and agent governance.}
\label{tab:survey_comparison}
\small
\setlength{\tabcolsep}{4pt}
\begin{tabular}{p{3.2cm}p{3.3cm}p{2.6cm}p{2.7cm}p{2.9cm}p{2.6cm}}
\toprule
\textbf{Work \& Reference} & \textbf{Primary Focus} & \textbf{Data Access \& IAM (RQ1)} & \textbf{Prompt Injection Defense (RQ2)} & \textbf{Tool \& Action Safety (RQ3)} & \textbf{Standards Alignment (RQ4)} \\
\midrule
\multicolumn{6}{l}{\textit{\textbf{Existing Surveys and Meta-Reviews}}} \\
\midrule
Dong et al.~\cite{dong2025safeguardingsurvey} [Survey] & LLM guardrails, output safety, and alignment & No (conversational context only) & Yes (input/output guardrail classifiers) & No (unconstrained text generation) & Partial (general AI safety principles) \\
\addlinespace
Gao et al.~\cite{gao2023ragsurvey} [Survey] & Retrieval-augmented generation paradigms & No (semantic similarity without IAM) & No (corpus assumed trusted) & No (retrieval-to-generation only) & No (algorithmic benchmark focus) \\
\addlinespace
Wang et al.~\cite{wang2024agents} [Survey] & Autonomous agent architectures \& planning & No (unrestricted agent memory/context) & No (unprotected planning loops) & Partial (action selection; no safety checks) & No (agent capability benchmarks) \\
\addlinespace
Dev et al.~\cite{dev2025threat} [Meta-Review] & AI threat modeling and risk taxonomy & Partial (identifies data privacy threats) & Partial (identifies injection attack vectors) & Limited (threat cataloging; no runtime guards) & High (MITRE ATLAS, OWASP, NIST AI RMF) \\
\midrule
\multicolumn{6}{l}{\textit{\textbf{Representative Primary Research Papers}}} \\
\midrule
Jeong \& Lee~\cite{jeonglee2025permissionawarerag} [Research Paper] & Permission-aware RAG access filtering & Yes (Native AWS/GCP/Keycloak IAM filtering) & No (focuses strictly on authorization) & No (retrieval pipeline only) & Partial (enterprise cloud IAM integration) \\
\addlinespace
Debenedetti et al.~\cite{debenedetti2026camel} [Research Paper] & Indirect prompt injection defense & No (assumes pre-filtered context) & Yes (CaMeL data-flow privilege separation) & Partial (capability checks before tool calls) & Partial (AgentDojo threat benchmark) \\
\addlinespace
Mansoor et al.~\cite{mansoor2026verifiedtoolcalls} / Doshi et al.~\cite{doshi2026verifiablysafetooluse} [Research Paper] & Tool safety, formal invariants, and idempotency & No (tool execution layer only) & No (assumes uncorrupted tool inputs) & Yes (Postcondition verifiers, ledgers, Alloy models) & Partial (STPA hazard analysis) \\
\addlinespace
Kulkarni \& Ramanathan~\cite{kulkarni2025grails} [Framework / Paper] & Embedding ethical safeguards in software SDLC & Partial (data privacy design patterns) & No (SDLC process focus) & Limited (service-level validation patterns) & High (Responsible AI guidelines) \\
\midrule
\textbf{This survey} [Survey] & \textbf{End-to-end enterprise agent safeguarding} & \textbf{Yes (Dynamic IAM filtering \& multi-tenancy)} & \textbf{Yes (Dual-boundary structural isolation)} & \textbf{Yes (State machines, human approval, ledgers)} & \textbf{High (OWASP Top 10, NIST AI RMF, MITRE ATLAS)} \\
\bottomrule
\end{tabular}
\end{table*}

\subsection{Empirical Case Study: An Enterprise Codebase Audit}
\label{sec:audit}
To evaluate how these theoretical controls map to real-world software, we analyze an empirical static code audit of \emph{Enable Agents}, a production-oriented, full-stack multi-agent platform. The platform orchestrates complex business pipelines (such as supplier qualification and sales outreach) using Flask, React, PostgreSQL with \texttt{pgvector}, Redis, Celery, and LangGraph state machines. Table~\ref{tab:audit} maps the audited platform mechanisms to surveyed safeguards.

\begin{table}[t]
\centering
\caption{Audited platform mechanisms mapped to surveyed safeguard layers.}
\label{tab:audit}
\small
\setlength{\tabcolsep}{3pt}
\begin{tabular}{p{2.2cm}p{3.15cm}p{2.2cm}}
\toprule
\textbf{Layer / Focus} & \textbf{Audit finding} & \textbf{Survey link} \\
\midrule
Retrieval (RQ1) & User filter depends on caller-specified user ID; no per-document ACLs & \cite{jeonglee2025permissionawarerag} \\
Injection (RQ2) & Retrieved text concatenated into document-chat prompts; no isolation found & \cite{greshake2023indirect,debenedetti2026camel} \\
Approval (RQ3) & Workflow interrupt for irreversible email; standalone send endpoint may bypass it & \cite{doshi2026verifiablysafetooluse,owasp2025llm} \\
Duplicate effects (RQ3) & Recipient/content ledger claim before send; 50-recipient cap & \cite{mansoor2026verifiedtoolcalls} \\
Cost/logging (RQ4) & Central gateway, budget reservations, and usage metadata; prompts/completions omitted & \cite{dev2025threat,kulkarni2025grails} \\
Evaluation (RQ4) & Keyword/JSON checks; no faithfulness metric & \cite{es2024ragas} \\
\bottomrule\end{tabular}
\end{table}

\textbf{Implemented Strengths.} The codebase strongly reflects the \emph{Determinism Principle}: critical boundaries are enforced through native Python logic rather than delegated to language models:
\begin{itemize}
  \item \emph{State-Machine Approval Checkpoints:} In LangGraph workflows (\texttt{graph.py}), stages marked with irreversible side effects (\texttt{side\_effect="irreversible"}, e.g., bulk email dispatch) trigger an explicit \texttt{interrupt()}, pausing execution until a human administrator approves, edits, or rejects the candidate payload via \texttt{POST /api/workflows/instances/:id/resume}.
  \item \emph{Durable Claim-Before-Send Action Ledgers:} To mitigate non-atomic tool failures \cite{mansoor2026verifiedtoolcalls}, \texttt{WorkflowSendLedger} computes a SHA-1 content hash of outbound emails, asserting a database claim prior to dispatch. If worker processes crash or retry, existing claims prevent duplicate execution ($n_k \le 1$).
  \item \emph{Multi-Tier Financial Guardrails:} The model gateway (\texttt{core/ai\_client.py}) interfaces with \texttt{core/budget.py} to compute pre-flight cost estimations and reserve funds in \texttt{budget\_reservations}, blocking execution with HTTP 402 if project or team spending limits are exceeded.
\end{itemize}

\textbf{Architectural Vulnerabilities and Gaps.} Conversely, static code analysis reveals critical security deficits common across industrial implementations:
\begin{enumerate}
  \item \emph{Coarse Multi-Tenancy Scoping:} In an ideal permission-aware architecture, authorized context $C(q,u)$ for query $q$ and user $u$ is filtered against each document source's native IAM policy:
  \[
  C(q,u) = \{d \in \operatorname{Top}_k(q) : \operatorname{IAM}_{s(d)}(u, d, \text{read}) = \text{allow}\}.
  \]
  In the audited platform, vector search applies an uploader filter only if a caller explicitly supplies a user identifier. Letting $D$ be the indexed corpus and $u=\bot$ denote an omitted identifier, the effective context behaves as $C'(q,u) = \operatorname{Top}_k(q; D_u)$, where:
  \[
  D_u = \begin{cases}
  D, & u = \bot, \\
  \{d : \operatorname{owner}(d) = u\}, & \text{otherwise.}
  \end{cases}
  \]
  This uploader-scoping model lacks fine-grained document ACLs and risks cross-tenant exposure if callers omit parameters.
  \item \emph{Absence of Prompt-Injection Isolation:} In the document question-answering service, retrieved chunk text is directly concatenated into prompt templates without XML delimiter framing or capability tracking. An adversarial instruction embedded within an ingested PDF could directly hijack the conversational agent, highlighting the necessity of structural defenses like CaMeL \cite{debenedetti2026camel}.
  \item \emph{Heuristic-Only Testing:} Continuous integration fixtures verify keyword presence and JSON keys, lacking semantic hallucination detection (e.g., RAGAs \cite{es2024ragas}) or adversarial robustness benchmarking.
\end{enumerate}

To evaluate these layers rigorously against policy ground truth, retrieval access control must be measured via False Authorization Rate ($\mathrm{FAR}$) and False Rejection Rate ($\mathrm{FRR}$):
\[
\mathrm{FAR} = \frac{\mathrm{FP}}{\mathrm{FP} + \mathrm{TN}}, \qquad \mathrm{FRR} = \frac{\mathrm{FN}}{\mathrm{FN} + \mathrm{TP}},
\]
where $\mathrm{FP}$ counts unauthorized chunks exposed, $\mathrm{TN}$ unauthorized chunks denied, $\mathrm{FN}$ authorized chunks withheld, and $\mathrm{TP}$ authorized chunks correctly retrieved. Tool safety is measured by postcondition verification rate and duplicate-effect rate $\mathcal{R}_{\mathrm{dup}} = N_{\mathrm{dup}} / N_{\mathrm{exec}}$, while threat defenses require reporting Attack Success Rate ($\mathrm{ASR}$) alongside Benign Refusal Rate ($\mathrm{BRR}$).

% ==================================================================
% 5. RESEARCH GAPS AND OPEN CHALLENGES
% ==================================================================
\section{Research Gaps and Open Challenges}
\label{sec:gaps}
Our synthesis of literature and codebase analysis reveals four critical open research challenges:
\begin{enumerate}
  \item \textbf{Dynamic Synchronization and Permission Revocation:} Current permission-aware RAG research (e.g., \cite{jeonglee2025permissionawarerag}) evaluates static IAM snapshots. In production enterprises, permissions, group memberships, and documents change continuously. Investigating zero-delay cache invalidation, deletion propagation (tombstones), and dynamic re-indexing without leaking revoked context remains an unaddressed challenge.
  \item \textbf{Production vs. Benchmark Evaluation Fidelity:} Existing security benchmarks (e.g., HotpotQA, AgentDojo) rely on synthetic, narrow test cases. Real enterprise environments feature complex multi-hop dependencies, ambiguous human organizational roles, and strict latency service level agreements (SLAs). Developing standardized enterprise agent benchmarks with realistic threat landscapes is urgently required.
  \item \textbf{The Security-Utility-Latency Frontier:} Structural defenses such as CaMeL \cite{debenedetti2026camel} achieve impressive security guarantees but incur token overhead and task degradation (-7\% completion rate). Similarly, query-time IAM checks introduce measurable latency. Future work must explore speculative permission evaluation, hardware acceleration, and lightweight verification models to balance safety with user experience.
  \item \textbf{Unified End-to-End Forensic Traceability:} While individual systems record token expenditure (\texttt{AIUsageLog}) or action claims (\texttt{WorkflowSendLedger}), enterprise compliance requires an immutable, privacy-preserving audit trail linking original source provenance, IAM evaluation results, model reasoning tokens, human approval signatures, and external API postconditions into a unified forensic chain.
\end{enumerate}

% ==================================================================
% 6. CONCLUSION
% ==================================================================
\section{Conclusion}
\label{sec:conclusion}
Autonomous enterprise AI agents represent a paradigm shift in workflow automation, yet their integration with corporate knowledge bases and external APIs creates systemic vulnerabilities across data privacy, model integrity, and operational execution. This survey has synthesized contemporary literature into a four-layer architectural taxonomy: permission-aware data governance, threat protection, governed tool execution, and operational observability. 

Our core finding is that enterprise safeguarding cannot be achieved through model-level alignment or prompt engineering alone. Rather, enterprise safeguarding is fundamentally a systems-engineering discipline grounded in the \emph{Determinism Principle}: deterministic software code must enforce access control before data enters model context, treat retrieved data as untrusted regardless of origin, checkpoint irreversible real-world actions via state machines, and independently verify tool postconditions. Future research must bridge the divide between theoretical security proofs and high-throughput, dynamic enterprise environments.

% ==================================================================
% REFERENCES
% ==================================================================
\bibliographystyle{ACM-Reference-Format}
\bibliography{enterprise_ai_recommended_sources}

\end{document}
